\clearpage
\section{Étape III \textendash{} Analyse stratégique}

L'analyse stratégique confronte ce que nous sommes à ce que nous subissons. Elle
se compose de deux diagnostics complémentaires : un diagnostic interne, qui
recense nos forces et nos faiblesses selon la grille des ressources d'Edith
Penrose, et un diagnostic externe, qui identifie les opportunités et les menaces
de notre environnement, au niveau micro-économique par le modèle des forces
concurrentielles de Michael Porter, puis au niveau macro-économique par le
modèle PESTEL.

\subsection{Hypothèses de travail retenues}

Le diagnostic interne repose sur les décisions suivantes, arrêtées collectivement
et qui engagent l'ensemble du business plan.

\begin{center}
\begin{tabularx}{\textwidth}{P{4cm}L}
\toprule
\thead{Décision} & \thead{Choix retenu} \\
\midrule
Équipe & Trois associés à parts égales, diplômés du BUT Informatique. \\
Capital social & 15 000 euros d'apport personnel en numéraire, sans emprunt au
démarrage. \\
Rémunération & Aucune rémunération des associés en année 1, qui conservent une
activité salariée en parallèle. \\
Hébergement & Domicile d'un associé en année 1, pépinière d'entreprises en
année 2. \\
Conception électronique & Externalisée auprès d'un bureau d'études spécialisé. \\
Fabrication & Assemblage en France chez un sous-traitant électronique, sans
outil de production en propre. \\
Produit & Réduction active de bruit sans Bluetooth ni application. \\
\bottomrule
\end{tabularx}
\end{center}

\begin{hypothese}
La répartition nominative des rôles entre les trois associés est présentée
ci-après à titre de proposition. Elle reste à arbitrer définitivement en
fonction des appétences de chacun.
\end{hypothese}

\subsection{Diagnostic interne}

\subsubsection{Organisation de l'équipe fondatrice}

La structure repose sur trois pôles de responsabilité, un par associé, afin
d'éviter la dilution des décisions et de clarifier qui répond de quoi.

\begin{center}
\begin{tabularx}{\textwidth}{P{3.4cm}P{3.2cm}L}
\toprule
\thead{Pôle} & \thead{Associé} & \thead{Périmètre} \\
\midrule
Produit et technique & \etudiantun & Cahier des charges, relation avec le bureau
d'études, suivi de fabrication, contrôle qualité. \\
Commercial et marketing & \etudiantdeux & Site de vente, communication,
acquisition de clients, service après-vente. \\
Administratif et financier & \etudianttrois & Juridique, comptabilité,
financement, conformité réglementaire et certification. \\
\bottomrule
\end{tabularx}
\end{center}

\subsubsection{Tableau de diagnostic interne}

\begin{center}
\begin{xltabular}{\textwidth}{P{3.2cm}LL}
\toprule
\thead{Ressource} & \thead{\textcolor{bpvert}{Forces}} & \thead{\textcolor{bprouge}{Faiblesses}} \\
\midrule
\endfirsthead
\toprule
\thead{Ressource} & \thead{\textcolor{bpvert}{Forces}} & \thead{\textcolor{bprouge}{Faiblesses}} \\
\midrule
\endhead

\textbf{Physiques}
&
Aucun local à financer en année 1, l'activité étant hébergée au domicile d'un
associé. Aucun outil de production à acquérir, la fabrication étant
sous-traitée. Aucun stock en année 1, puisque la commercialisation ne débute
qu'en année 2. Implantation dans une agglomération bien desservie, facilitant
l'accès aux réseaux d'accompagnement.
&
Absence totale d'atelier et de moyens de mesure acoustique, ce qui interdit de
valider un prototype en interne. Dépendance complète au sous-traitant pour la
fabrication. Absence de lieu d'accueil professionnel, susceptible de peser sur
la crédibilité lors des premiers contacts fournisseurs. \\

\textbf{Financières}
&
Capital social de 15 000 euros entièrement libéré, sans dette financière au
démarrage, donc aucune charge d'intérêts en année 1. Absence de rémunération
des associés, ce qui supprime le premier poste de charges d'un compte de
résultat. Structure de coûts légère et essentiellement variable.
&
Trésorerie insuffisante pour financer seule l'outillage et la première série,
ce qui rendra un emprunt indispensable en année 2. Aucun revenu en année 1,
donc un exercice nécessairement déficitaire. Aucune capacité à absorber un
dépassement important sur le développement. \\

\textbf{Humaines}
&
Trois associés aux rôles clairement répartis, ce qui évite la dispersion.
Complémentarité entre un profil technique, un profil commercial et un profil
administratif. Prise de décision rapide compte tenu de la taille de l'équipe.
Motivation forte et engagement de long terme.
&
Aucune compétence en électronique ni en acoustique au sein de l'équipe, alors
qu'il s'agit du c\oe{}ur du produit. Aucune expérience entrepreneuriale
antérieure. Disponibilité limitée en année 1, les associés conservant une
activité salariée. Absence de réseau établi dans l'industrie du son. \\

\textbf{Technologiques}
&
Maîtrise du développement web et du commerce en ligne, ce qui permet de
construire et d'exploiter le canal de vente sans prestataire. Aucune
dépendance logicielle, le produit ne comportant ni application ni connexion.
Absence de traitement de données personnelles, ce qui écarte les obligations
liées au RGPD.
&
Aucun brevet ni savoir-faire propriétaire, donc aucune barrière à l'imitation.
Toute la valeur technique du produit repose sur un prestataire externe.
Dépendance aux composants importés, notamment les processeurs de traitement du
signal et les batteries. \\

\textbf{Organisationnelles}
&
Structure très légère, sans hiérarchie ni processus lourds, donc réactive.
Trois pôles identifiés, chacun avec un responsable désigné. Décisions
collégiales sur les arbitrages engageant l'entreprise.
&
Gouvernance non formalisée et absence de procédures écrites. Risque de blocage
en cas de désaccord entre trois associés à parts égales. Aucune procédure de
continuité en cas d'indisponibilité d'un associé. \\

\textbf{Mercatiques et commerciales}
&
Positionnement clair et différencié sur le confort, l'autonomie réelle et la
simplicité. Argument d'assemblage en France, dans un contexte où l'origine
française est un critère d'achat reconnu. Vente directe sans intermédiaire,
donc marge préservée et relation client maîtrisée. Marché international
accessible sans investissement supplémentaire.
&
Notoriété nulle au lancement, face à des marques déjà installées. Aucun fichier
client et aucune communauté. Budget de communication très limité. Absence
d'historique d'avis clients, alors que ceux-ci sont déterminants dans l'achat
en ligne d'un produit technique. \\
\bottomrule
\end{xltabular}
\end{center}

\clearpage
\subsection{Diagnostic externe \textendash{} micro-environnement}

Le micro-environnement est analysé selon le modèle des cinq forces
concurrentielles de Michael Porter, complété par la sixième force que
constituent les contraintes réglementaires.

\subsubsection{Analyse de la concurrence}

Le marché des dispositifs d'aide au sommeil par le son se structure autour de
deux approches technologiques distinctes. La première supprime le bruit par
réduction active, la seconde le masque en diffusant un son de couverture. Cette
distinction est essentielle car elle sépare nos concurrents directs de nos
concurrents indirects.

\begin{center}
\begin{xltabular}{\textwidth}{P{2.8cm}P{2.4cm}P{2.2cm}L}
\toprule
\thead{Concurrent} & \thead{Approche} & \thead{Prix indicatif} & \thead{Caractéristiques et limites} \\
\midrule
\endfirsthead
\toprule
\thead{Concurrent} & \thead{Approche} & \thead{Prix indicatif} & \thead{Caractéristiques et limites} \\
\midrule
\endhead

QuietOn 4
&
Réduction active
&
259 euros
&
Concurrent direct le plus proche de notre positionnement. Sans Bluetooth ni
application, autonomie annoncée jusqu'à 28 heures. Prix élevé et absence de
fonction audio relevés comme principaux reproches. Un test indépendant signale
une difficulté à obtenir un maintien confortable (Le Café du Geek, 2025 ;
CNN Underscored, 2026). \\

Soundcore Sleep A30
&
Réduction active et masquage
&
À partir de 196 euros
&
Adossé au groupe Anker, donc doté d'une force de frappe industrielle et
marketing supérieure. Produit plus complet, avec suivi du sommeil et diffusion
audio, mais autonomie ramenée à environ six heures lorsque la réduction active
est engagée (Idealo, 2026). \\

Ozlo Sleepbuds
&
Masquage sonore
&
Gamme haute
&
Concurrent indirect. Ne supprime pas le bruit mais le couvre. Confirme
l'existence d'une approche alternative au même problème. \\

Bouchons en mousse
&
Atténuation passive
&
Moins de 5 euros
&
Concurrent indirect le plus redoutable par son prix. Atténue mal les basses
fréquences, ce qui constitue précisément notre terrain de différenciation. \\
\bottomrule
\end{xltabular}
\end{center}

\begin{hypothese}
Les chiffres d'affaires et parts de marché de ces acteurs ne sont pas publiés
de manière exploitable pour un marché de niche. Les prix indiqués sont relevés
sur les sites officiels et comparateurs, et devront être réactualisés avant le
rendu final.
\end{hypothese}

\subsubsection{Tableau de diagnostic externe, micro-environnement}

\begin{center}
\begin{xltabular}{\textwidth}{P{3.2cm}LL}
\toprule
\thead{Force concurrentielle} & \thead{\textcolor{bpvert}{Opportunités}} & \thead{\textcolor{bprouge}{Menaces}} \\
\midrule
\endfirsthead
\toprule
\thead{Force concurrentielle} & \thead{\textcolor{bpvert}{Opportunités}} & \thead{\textcolor{bprouge}{Menaces}} \\
\midrule
\endhead

\textbf{Intensité concurrentielle}
&
Marché de niche encore restreint, où aucun acteur ne domine réellement. Aucun
concurrent ne résout simultanément le confort, l'autonomie et la simplicité.
Segment récent, laissant place à un positionnement nouveau.
&
Concurrents déjà installés et mieux financés, disposant d'une notoriété et
d'avis clients accumulés. Comparaison directe et immédiate par le consommateur
sur les places de marché. \\

\textbf{Pouvoir de négociation des clients}
&
Clientèle de particuliers dispersée, donc sans pouvoir de négociation
individuel sur les prix. Disposition à payer démontrée par les tarifs
pratiqués sur le marché.
&
Coût de changement nul : le client compare en quelques minutes et achète
ailleurs. Les avis en ligne conditionnent fortement la décision d'achat, ce qui
pénalise une marque sans historique. Exigence de résultat immédiate sur un
produit vendu cher. \\

\textbf{Pouvoir de négociation des fournisseurs}
&
Existence d'un tissu de sous-traitants électroniques en France, compatible avec
notre choix d'assemblage national. Possibilité de mettre plusieurs bureaux
d'études en concurrence lors de la phase de conception.
&
Dépendance forte à un nombre réduit de partenaires pour la conception et
l'assemblage. Volumes faibles au démarrage, donc faible pouvoir de négociation
et quantités minimales de commande imposées. Approvisionnement en composants
critiques, processeurs et batteries, concentré hors d'Europe. \\

\textbf{Menace des nouveaux entrants}
&
Barrière réelle liée au savoir-faire acoustique et à la miniaturisation, qui
freine les entrants opportunistes.
&
Un acteur établi de l'audio peut entrer sur le segment à tout moment avec des
moyens sans commune mesure. Aucune protection par brevet de notre côté.
Possibilité de copie rapide par des fabricants à bas coût. \\

\textbf{Menace des produits de substitution}
&
Aucune substitution ne traite efficacement les basses fréquences, ce qui
préserve la pertinence de la réduction active.
&
Menace très forte et multiforme : bouchons en mousse à quelques euros,
générateurs de bruit blanc, applications gratuites, isolation phonique du
logement, voire traitement médical du ronflement du conjoint. Le rapport entre
notre prix et celui des substituts est considérable. \\

\textbf{Contraintes réglementaires}
&
Les exigences de certification constituent aussi une barrière à l'entrée pour
les concurrents non européens mal préparés.
&
Marquage CE obligatoire, réglementation sur la sécurité et le transport des
batteries lithium, filière de traitement des déchets électroniques, garantie
légale de conformité. Contrôle de la DGCCRF sur les allégations commerciales,
notamment sur l'origine française et sur toute allégation de santé. \\
\bottomrule
\end{xltabular}
\end{center}

\clearpage
\subsection{Diagnostic externe \textendash{} macro-environnement}

Le macro-environnement est analysé selon le modèle PESTEL, qui couvre les six
dimensions de l'environnement général de l'entreprise.

\begin{center}
\begin{xltabular}{\textwidth}{P{3.2cm}LL}
\toprule
\thead{Environnement} & \thead{\textcolor{bpvert}{Opportunités}} & \thead{\textcolor{bprouge}{Menaces}} \\
\midrule
\endfirsthead
\toprule
\thead{Environnement} & \thead{\textcolor{bpvert}{Opportunités}} & \thead{\textcolor{bprouge}{Menaces}} \\
\midrule
\endhead

\textbf{Politique}
&
Politiques publiques de soutien à la réindustrialisation et à la relocalisation
de production, cohérentes avec notre choix d'assemblage en France. Dispositifs
d'accompagnement à la création d'entreprise accessibles aux jeunes diplômés.
&
Instabilité des dispositifs d'aide, dont les conditions évoluent d'une année
sur l'autre. Tensions commerciales internationales susceptibles d'affecter
l'approvisionnement en composants. \\

\textbf{Économique}
&
Croissance du marché des produits liés au bien-être et au sommeil. Prix élevés
déjà acceptés sur ce segment, ce qui autorise une marge confortable. Vente en
ligne directe, sans coût de réseau physique.
&
Sensibilité d'un produit non essentiel aux arbitrages de pouvoir d'achat.
Inflation sur les composants électroniques. Coût de production en France
nettement supérieur à celui d'un assemblage en Asie, ce qui contraint le prix
de vente. \\

\textbf{Socioculturel}
&
Le sommeil devient une préoccupation de santé majeure : 42 \% des Français
placent le fait de bien dormir en tête des piliers de la santé, et 36 \% se
déclarent gênés par le bruit la nuit (INSV, 2026). L'origine française est un
critère d'achat largement partagé, 85 \% des Français déclarant acheter des
produits fabriqués en France. Sensibilité croissante à la protection des
données personnelles, favorable à un produit sans application.
&
Débat médiatique naissant sur les effets d'un usage prolongé de la réduction
active de bruit sur le traitement auditif, susceptible de peser sur l'image de
la technologie. Méfiance d'une partie des consommateurs envers les jeunes
marques inconnues sur un produit technique coûteux. \\

\textbf{Technologique}
&
Progrès continus sur la miniaturisation, les processeurs de traitement du
signal et la densité des batteries, qui rendent le produit plus réalisable
qu'il y a quelques années. Maturité des outils de vente en ligne.
&
Obsolescence rapide des composants électroniques. Risque qu'un concurrent
franchisse une étape technologique décisive sur le confort ou l'autonomie.
Difficulté technique persistante du traitement du ronflement, son non
stationnaire. \\

\textbf{Écologique}
&
Attente forte de durabilité et de réparabilité, à laquelle un produit simple,
sans application ni obsolescence logicielle, répond naturellement. Assemblage
en France réduisant l'empreinte du transport.
&
Obligations liées aux déchets d'équipements électriques et électroniques, avec
éco-contribution et filière de reprise. Contraintes croissantes sur le
recyclage et la fin de vie des batteries. Critique de principe visant tout
produit électronique jetable. \\

\textbf{Légal}
&
L'absence d'application et de collecte de données écarte les obligations liées
au RGPD, ce qui simplifie notablement notre conformité.
&
Marquage CE et normes de sécurité électrique obligatoires. Réglementation
stricte sur l'usage de la mention d'origine française, contrôlée par la DGCCRF
et la douane. Interdiction de toute allégation de santé, sous peine de
basculer dans le régime des dispositifs médicaux. Garantie légale de
conformité et droit de rétractation applicables à la vente à distance. \\
\bottomrule
\end{xltabular}
\end{center}

\clearpage
\subsection{Synthèse de l'analyse stratégique}

\subsubsection{Matrice SWOT consolidée}

\begin{center}
\begin{tabularx}{\textwidth}{LL}
\toprule
\thead{\textcolor{bpvert}{Forces}} & \thead{\textcolor{bprouge}{Faiblesses}} \\
\midrule
Capital sans dette, structure de coûts légère, aucun local ni outil de
production à financer. Équipe complémentaire aux rôles répartis. Positionnement
différencié et documenté. Vente directe sans intermédiaire. Aucune donnée
personnelle traitée.
&
Aucune compétence en électronique ni en acoustique. Notoriété nulle et absence
d'avis clients. Trésorerie insuffisante pour la phase de série. Aucun brevet.
Disponibilité limitée en année 1. \\
\midrule
\thead{\textcolor{bpvert}{Opportunités}} & \thead{\textcolor{bprouge}{Menaces}} \\
\midrule
Préoccupation croissante pour le sommeil et forte gêne déclarée face au bruit.
Marché de niche sans acteur dominant. Attrait avéré pour l'origine française.
Aucun concurrent ne résout à la fois le confort, l'autonomie et la simplicité.
Sensibilité à la protection des données.
&
Substituts à très bas prix. Entrée possible d'un acteur majeur de l'audio.
Dépendance aux fournisseurs et aux composants importés. Surcoût de la
fabrication française. Réglementation lourde sur la certification et les
allégations. Débat public sur les effets de la réduction active de bruit. \\
\bottomrule
\end{tabularx}
\end{center}

\subsubsection{Enseignements stratégiques}

Trois conclusions se dégagent de cette analyse et orienteront les étapes
suivantes.

\textbf{La faiblesse la plus critique est technique, pas financière.} L'absence
de compétence en électronique et en acoustique porte sur le c\oe{}ur même du
produit. Le choix du bureau d'études externe n'est donc pas un détail
d'organisation mais la décision la plus structurante du projet. Elle sera
traitée en priorité à l'étape VIII.

\textbf{La menace la plus sérieuse vient des substituts, pas des concurrents
directs.} Un bouchon en mousse coûte moins de cinq euros. L'enjeu commercial
n'est donc pas de convaincre le client de nous préférer à QuietOn, mais de le
convaincre qu'une solution à deux cents euros se justifie. C'est ce point que
l'étude de marché de l'étape V devra trancher en priorité.

\textbf{L'opportunité la plus exploitable est le croisement entre le confort et
l'origine française.} Le confort est le défaut relevé chez les concurrents, et
l'origine française est un critère d'achat partagé par une large majorité de
consommateurs. Ces deux éléments fondent le positionnement du plan de
marchéage développé à l'étape IV.
